\chapter{Paradoxos de Autorreferência e Definibilidade}
\label{chap:apendiceAutorreferenciaDefinibilidade}

Uma linha une todos os paradoxos deste apêndice: a autorreferência. Quando uma linguagem ou sistema é expressivo o suficiente para falar sobre si mesmo, surgem inevitavelmente afirmações que subvertem sua própria lógica. Veremos que a mesma técnica da diagonal utilizada por Cantor para provar a não enumerabilidade de \(\mathbb{R}\) no \autoref{chap:infinitos} reaparece, sob diferentes formas, em cada um dos paradoxos a seguir.

\section{O Barbeiro de Russell (versão popular)}

Apresentar a versão informal: o barbeiro que barbeia exatamente quem
não se barbeia. Usar como porta de entrada acessível antes da versão
formal, alinhado com o objetivo declarado de escrever para graduandos.


\section{O Paradoxo de Russell (versão formal, 1901)}

Retomar o conjunto \(B = \{x : x \not\in x\}\) já apresentado no Capítulo 2,
agora sob a perspectiva histórica e suas consequências para os
fundamentos da matemática. Discutir a resposta de Zermelo via
Axioma de Separação (Axioma 6 do trabalho).


\section{O Paradoxo de Richard (1905)}

Apresentar a construção de Richard: lista de todas as descrições finitas
de números reais em língua francesa, seguida de diagonalização.
Conectar explicitamente com o Argumento da Diagonal do Capítulo 4.


\section{O Paradoxo de Berry (1906)}

Apresentar: "o menor inteiro positivo que não pode ser descrito com
menos de quinze palavras", frase com 14 palavras.
Conectar com o Exemplo do conjunto A (p. 7 do trabalho), do qual
este paradoxo é uma versão simplificada e ainda mais direta.
Mencionar a formalização moderna via complexidade de Kolmogorov.


\section{O Paradoxo de Grelling-Nelson (1908)}

Apresentar adjetivos autológicos e heterológicos.
Mostrar que o problema da autorreferência não é exclusivo da teoria
dos conjuntos, mas emerge em qualquer sistema suficientemente expressivo.

\section{O Paradoxo de Curry (1942)}

Apresentar a sentença autorreferente \(C_Q\) que, sem uso de negação,
deriva qualquer proposição \(Q\) arbitrária via \textit{modus ponens}.
Conectar com o Axioma de Separação (Axioma~6): mostrar que a
variante conjuntista do paradoxo, \(C_Q = \{x \mid x \in x
\rightarrow Q\}\), é bloqueada pela exigência de separar a partir
de um conjunto preexistente. Discutir o fenômeno de explosão
lógica (\textit{ex contradictione quodlibet}) e contrastar com
Russell: enquanto Russell exige negação, Curry exige apenas
implicação, tornando o colapso lógico ainda mais geral.

\section{O Teorema da Incompletude de Gödel (1931)}

Apresentar como o mesmo método diagonal de Cantor é usado para
construir uma afirmação que diz ``esta afirmação não é demonstrável''.
Enfatizar a conexão técnica com o Capítulo 4: diagonalização como
ferramenta unificadora. Tratar com cuidado a distinção entre
paradoxo e teorema.